\documentclass[11pt]{article}
\usepackage[T1]{fontenc}
\usepackage{lmodern,amsmath,amssymb,amsthm,mathtools}
\usepackage[margin=1.05in]{geometry}
\usepackage{longtable,array}
\usepackage[hidelinks]{hyperref}
\newtheorem{theorem}{Theorem}[section]
\newtheorem{lemma}[theorem]{Lemma}
\newtheorem{proposition}[theorem]{Proposition}
\newtheorem{corollary}[theorem]{Corollary}
\theoremstyle{remark}\newtheorem{remark}[theorem]{Remark}
\newcommand{\R}{\mathbb R}\newcommand{\C}{\mathbb C}
\newcommand{\N}{\mathbb N}
\newcommand{\abs}[1]{\left|#1\right|}
\newcommand{\norm}[1]{\left\lVert#1\right\rVert}
\newcommand{\F}{F_{a,t}}
\newcommand{\code}[1]{{\small\nolinkurl{#1}}}
\title{Endpoint singularities and nonreal zeros\protect\\
in fractional theta deformations}
\author{Project Stardust}
\date{Reviewed editorial draft, October 29, 2026}
\begin{document}
\maketitle
\begin{abstract}
We study the cosine transforms of the arithmetic theta kernel multiplied
by $\exp(tu^{2a})$ on the positive half-line. For $1<a<2$ and $t\neq0$,
the transform is entire but has only finitely many real zeros and
infinitely many nonreal zeros. The mechanism is the endpoint asymptotic
$F_{a,t}(x)\sim tK(0)\Gamma(2a+1)\sin(\pi(a-1))x^{-2a-1}$.
Near quadratic order a uniform absolute error gives a real-zero radius bounded by a constant times the inverse squared order gap, while full rectangular relative control fails. We give the canceled-derivative proof on both sides of the cubic
endpoint and treat that endpoint separately. The transforms converge
locally uniformly in the complex plane and uniformly on the entire
real axis as the order varies. Nevertheless their global zero sets
need not be real. At a limiting order with only real simple zeros,
every selected nonreal-zero sequence escapes compact sets. The exact
domain and quantifiers agree with a frozen Lean development. The
Fourier endpoint mechanism has classical antecedents; novelty of the
arithmetic specialization and its combination with order limits is
not asserted.
\end{abstract}

\section{The family and the main obstruction}
For $u\geq0$ set
\begin{equation}\label{eq:kernel}
 K(u)=\sum_{n=1}^{\infty}
 \left(2\pi^2n^4e^{9u}-3\pi n^2e^{5u}\right)
 e^{-\pi n^2e^{4u}}.
\end{equation}
The kernel has an even analytic continuation near the real axis,
whose restriction to negative real arguments we also denote by $K$.
For $a>0$, $t\in\R$, and $z\in\C$, define
\begin{equation}\label{eq:flow}
 F_{a,t}(z)=\int_0^\infty e^{tu^{2a}}K(u)\cos(zu)\,du.
\end{equation}
The time sign is positive. Rapid kernel decay makes the integral
meaningful for every real time. At integer $a=k$, the tilt is
$\exp(tu^{2k})$; thus $a=2$ is the quartic order, $a=3$ is the sixth
order, and $a=3/2$ is the cubic endpoint between the ranges considered
below. These parameter conventions will not be interchanged.

\begin{theorem}\label{thm:main}
Fix $1<a<2$ and $t\neq0$. The function $F_{a,t}$ is a nonzero even
entire function, real on the real axis. As $x\to+\infty$,
\begin{equation}\label{eq:asymptotic}
 x^{2a+1}F_{a,t}(x)\longrightarrow
 L(a,t):=tK(0)\Gamma(2a+1)\sin(\pi(a-1))\neq0.
\end{equation}
Its real zero set is finite, and its nonreal zero set is infinite.
Consequently there is no real $T$ such that all zeros of $F_{a,t}$
are real for every $t\geq T$.
\end{theorem}

\begin{theorem}\label{thm:limit}
For fixed $t$ and $a_0>0$, $F_{a,t}\to F_{a_0,t}$ locally uniformly
on $\C$ and uniformly on $\R$ as $a\to a_0$. If all zeros of
$F_{a_0,t}$ are real and simple, and
\[
 a_n\to a_0,\qquad F_{a_n,t}(z_n)=0,\qquad \operatorname{Im}z_n\neq0,
\]
then $|z_n|\to\infty$.
For each fixed $t\neq0$, there is an escaping selection of nonreal
zeros along $a_n=1+1/(n+3)$ without any assumption about the zeros
of $F_{1,t}$.
\end{theorem}

The two escape statements have different quantifiers. The last one
chooses a sufficiently distant zero separately at each order.
The preceding statement applies to \emph{every} selection, and
requires both reality and simplicity of the limiting zeros. No
reality or simplicity hypothesis at quadratic order is silently
deduced from convergence.

P\'olya's work on entire Fourier transforms is the natural historical
context~\cite{polya}. The non-even-power examples identified in the
literature already warn that endpoint regularity can affect global
zero reality. Cardon's real-zero-preserving transform
theorems~\cite{cardon} provide another comparison: their hypotheses
are not replaced by positivity of the kernel in \eqref{eq:flow}.
The results here concern this precise arithmetic deformation and its
formal verification. A bounded preliminary search has not settled
priority for that specialization. Historical context is not used as
a premise of the analytic proofs below.

\section{Kernel identities and bounds}
We first record why the kernel has the properties needed at both
ends of the integral. They serve different purposes: evenness cancels
odd derivatives at zero, whereas rapid decay controls every tilt
and differentiation at infinity.

\begin{lemma}\label{lem:kernel}
The kernel is positive on $[0,\infty)$, is even and real analytic
on the real line, and for every integer $j\geq0$ satisfies a bound
\begin{equation}\label{eq:decay}
 |K^{(j)}(u)|\leq C_j\exp\left(C_j u-\frac\pi2 e^{4u}\right)
 \qquad(u\geq0)
\end{equation}
for a finite positive $C_j$. In particular $K(0)>0$ and every odd
derivative of $K$ vanishes at zero.
\end{lemma}
\begin{proof}
Every summand in \eqref{eq:kernel} is positive on $u\geq0$:
factoring out $\pi n^2e^{5u}$ leaves $2\pi n^2e^{4u}-3>0$.
The defining series converges absolutely. After any fixed number
of derivatives, each summand is a finite sum of terms bounded by
\[
 C n^M e^{Mu}e^{-\pi n^2e^{4u}}
\]
for some integer $M$ depending on the derivative order. With
$s=e^{4u}\geq1$, the inequality
$n^2s\geq(n^2+s)/2$ shows that their sum is at most
\[
 C e^{Mu-\pi s/2}\sum_{n\geq1}n^M e^{-\pi n^2/2}.
\]
The numerical series is finite. This proves \eqref{eq:decay} after
enlarging the constant and justifies termwise differentiation on
compact intervals.

For evenness, let $\vartheta(v)=\sum_{n\in\mathbb Z}e^{-\pi n^2v}$
and $Q(u)=e^u\vartheta(e^{4u})$. The theta transformation
$\vartheta(v)=v^{-1/2}\vartheta(1/v)$ gives $Q(-u)=Q(u)$.
Differentiating the absolutely convergent series shows
\begin{equation}\label{eq:theta}
 K(u)=\frac1{16}\big(Q''(u)-Q(u)\big).
\end{equation}
Hence $K$ is even. The theta transformation itself follows from
Poisson summation of the Gaussian: periodize $e^{-\pi v x^2}$,
whose $n$th Fourier coefficient is $v^{-1/2}e^{-\pi n^2/v}$,
and evaluate the absolutely convergent Fourier series at zero.
Absolute and locally uniform convergence justify this computation.
The same defining theta series is holomorphic where
$\operatorname{Re}(e^{4u})>0$, in particular on the strip
$|\operatorname{Im}u|<\pi/8$. It follows that $K$ is analytic near
the real line. Evenness gives the odd-derivative cancellations.
\end{proof}

The Gaussian Fourier coefficient in this proof follows by scaling
the real Gaussian integral and differentiating its Fourier transform:
integration by parts gives $\widehat g'(\xi)=-2\pi\xi\widehat g(\xi)/v$,
and $\widehat g(0)=v^{-1/2}$. This determines the coefficient and
fixes the normalization in \eqref{eq:theta}.

\begin{lemma}\label{lem:entire}
For every $a>0$ and $t\in\R$, the integral \eqref{eq:flow} is entire,
even, obeys $F_{a,t}(\bar z)=\overline{F_{a,t}(z)}$, and has
$F_{a,t}(0)>0$. Moreover, for some $C=C(a,t)>0$,
\begin{equation}\label{eq:growth}
 |F_{a,t}(z)|\leq C\exp(|z|^{3/2})\qquad(z\in\C).
\end{equation}
\end{lemma}
\begin{proof}
For $|z|\leq R$, the $j$th spatial derivative of the cosine is
bounded by $u^j e^{Ru}$. Bound \eqref{eq:decay} makes
$u^j\exp(tu^{2a}+Ru)K(u)$ absolutely integrable for every $j$.
Differentiation under the integral therefore gives entireness.
Evenness and conjugation symmetry follow from cosine, and positivity
at zero follows from the positive density.

For the growth bound put $r=|z|$. The elementary inequality
$ru\leq r^{3/2}+u^3$ holds for $r,u\geq0$: if $u\leq\sqrt r$
use the first term, and otherwise use the second. Thus
\[
 |F_{a,t}(z)|\leq e^{r^{3/2}}
 \int_0^\infty \exp(tu^{2a}+u^3)K(u)\,du.
\]
The integral is a finite positive constant by the same decay bound.
\end{proof}

The exponent $3/2$ is a convenient subquadratic bound, not an assertion
of the exact entire order. Its strict inequality below two will
be enough for the finite-zero rigidity argument.

\section{The endpoint Fourier calculation}
Write $\beta=2a\in(2,4)$ and
$g(u)=e^{tu^\beta}K(u)$ on $u>0$. All derivatives used below are
ordinary derivatives of this actual positive-axis function. We do
not differentiate a truncated formal jet beyond the orders for
which it represents those derivatives.

For $b>0$, the Laplace integral gives
\begin{equation}\label{eq:gamma}
 \int_0^\infty u^{b-1}e^{-u}e^{ixu}\,du
 =\Gamma(b)(1-ix)^{-b}.
\end{equation}
The right side uses the branch with $\arg(1-ix)\in(-\pi/2,\pi/2)$.
Indeed, for positive real $w$ the identity
$\int u^{b-1}e^{-wu}\,du=\Gamma(b)w^{-b}$ follows by substitution.
Both sides are holomorphic for $\operatorname{Re}w>0$, by dominated
differentiation on compact subsets, so the identity theorem extends
it to that half-plane and hence to $w=1-ix$.
For $x\to+\infty$, taking imaginary and real parts yields
\begin{align}
 x^b\int_0^\infty u^{b-1}e^{-u}\sin(xu)\,du
 &\longrightarrow\Gamma(b)\sin(\pi b/2),\label{eq:gammas}\\
 x^b\int_0^\infty u^{b-1}e^{-u}\cos(xu)\,du
 &\longrightarrow\Gamma(b)\cos(\pi b/2).\label{eq:gammac}
\end{align}
These formulas replace a heuristic oscillatory integral of a bare
power by an absolutely convergent integral with a known transform.
The subtraction of the exponential cutoff is what makes the actual
remainder differentiable with an integrable derivative.

\subsection{Cancellation estimates near zero}
Use the falling factorial $(\beta)_m=\beta(\beta-1)\cdots(\beta-m+1)$.
For $m=3$ or $4$, put
\begin{equation}\label{eq:remainder}
 C_m=tK(0)(\beta)_m,\qquad
 R_m(u)=g^{(m)}(u)-C_m u^{\beta-m}e^{-u}.
\end{equation}
The following estimates are for fixed $\beta,t$; no uniformity in
the order parameter is being asserted.

\begin{lemma}\label{lem:cancel}
If $2<\beta<3$, then $R_3,R_3'\in L^1(0,\infty)$,
$R_3(0+)=0$, and $R_3(u)\to0$ as $u\to\infty$.
If $3<\beta<4$, the same integrability and decay statements hold for
$R_4,R_4'$, with $R_4(0+)=K^{(4)}(0)$ in place of zero.
\end{lemma}
\begin{proof}
Near zero, expand only the regular factors:
\[
 K(u)=K(0)+\frac{K''(0)}2u^2+O(u^4),
 \qquad e^{tu^\beta}=1+tu^\beta+O(u^{2\beta}).
\]
The estimates persist after any of the finitely many positive-axis
derivatives required here. One justification is to write the
exponential remainder as $u^{2\beta}h(u^\beta)$ with $h$ analytic
near zero, and the kernel remainders as Taylor remainders with smooth
factors. Repeated product differentiation then gives, for $m=3,4$,
\begin{equation}\label{eq:expansion}
 g^{(m)}(u)=K^{(m)}(u)+C_m u^{\beta-m}
 +O(u^{\beta+2-m})+O(u^{2\beta-m}).
\end{equation}
Its next derivative has the differentiated power bounds, by that
same factorization argument, rather than by differentiating an
unqualified big-$O$ statement.

For $m=3$ the singular part left after the cutoff is
$C_3u^{\beta-3}(1-e^{-u})=O(u^{\beta-2})$.
The kernel term $K'''(u)$ is $O(u)$ by evenness.
Thus $R_3(0+)=0$, and locally
\[
 R_3'(u)=O(1)+O(u^{\beta-3})+O(u^{2\beta-4}).
\]
All exponents exceed $-1$ when $\beta>2$.

For $m=4$ the leftover cutoff term is $O(u^{\beta-3})$,
and the kernel term tends to $K^{(4)}(0)$. Its derivative satisfies
\[
 R_4'(u)=O(1)+O(u^{\beta-4})+O(u^{2\beta-5}),
\]
whose exponents exceed $-1$ when $\beta>3$.
The remainders themselves are locally integrable as well.
At infinity, repeated product differentiation gives finite sums of
powers of $u$ times $e^{tu^\beta}K^{(j)}(u)$; all are integrable
and tend to zero by \eqref{eq:decay}. The cutoff and its derivative
have the same properties at infinity.
\end{proof}

In particular the fourth-order remainder need not vanish at zero.
The forthcoming cosine integration by parts uses a sine boundary
factor, which kills its finite endpoint value. Assuming a zero
endpoint for that remainder would be an unnecessary and unjustified
strengthening.

\subsection{The range $2<\beta<3$}
Since $g'(0+)=0$, three integrations by parts give
\begin{equation}\label{eq:third}
 x^3F_{a,t}(x)=\int_0^\infty g'''(u)\sin(xu)\,du
 \qquad(x>0).
\end{equation}
For completeness, the first step is
$F=-x^{-1}\int g'\sin(xu)$; the second produces
$\int g'\sin(xu)=x^{-1}\int g''\cos(xu)$ because $g'(0+)=0$;
the third gives $\int g''\cos(xu)=-x^{-1}\int g'''\sin(xu)$.
All derivatives through the third order are integrable. Their
positive-axis endpoint limits and rapid decay justify the boundary
terms by first integrating over $[\varepsilon,M]$ and passing to
the limits.

By Lemma~\ref{lem:cancel}, another integration by parts gives
\[
 x\int_0^\infty R_3(u)\sin(xu)\,du
 =\int_0^\infty R_3'(u)\cos(xu)\,du.
\]
The right side is bounded by $\norm{R_3'}_{L^1}$; here and below
$\norm{\cdot}_{L^1}$ denotes the integral of the absolute value.
Set $b=\beta-2\in(0,1)$. Combining \eqref{eq:third} and
\eqref{eq:remainder}, multiplying by $x^b$, and using
\eqref{eq:gammas}, we obtain
\[
 x^{\beta+1}F_{a,t}(x)\longrightarrow
 C_3\Gamma(\beta-2)\sin(\pi(\beta-2)/2).
\]
The remainder vanishes because $x^bO(x^{-1})\to0$.
The recurrence for Gamma gives
$(\beta)_3\Gamma(\beta-2)=\Gamma(\beta+1)$, and
$\sin(\pi(\beta-2)/2)=\sin(\pi(a-1))$. This proves
\eqref{eq:asymptotic} in the first range.

\subsection{The range $3<\beta<4$}
Here $g'(0+)=g'''(0+)=0$. The previous integrations and one more
give
\begin{equation}\label{eq:fourth}
 x^4F_{a,t}(x)=\int_0^\infty g^{(4)}(u)\cos(xu)\,du.
\end{equation}
Lemma~\ref{lem:cancel} now yields
\[
 x\int_0^\infty R_4(u)\cos(xu)\,du
 =-\int_0^\infty R_4'(u)\sin(xu)\,du.
\]
At zero the boundary term is $R_4(0+)\sin0=0$; no vanishing of
$R_4(0+)$ is required. Set $b=\beta-3\in(0,1)$ and use
\eqref{eq:gammac}. The limit in \eqref{eq:fourth} becomes
\[
 C_4\Gamma(\beta-3)\cos(\pi(\beta-3)/2)
 =tK(0)\Gamma(\beta+1)\sin(\pi(a-1)).
\]
Again the remainder tends to zero after scaling because $b<1$.
This proves \eqref{eq:asymptotic} in the second range.

\subsection{The cubic endpoint $\beta=3$}
The positive-axis density is now $g(u)=e^{tu^3}K(u)$,
smooth up to zero on that half-line. Evenness of $K$ gives
$g'(0)=0$ and $g'''(0)=6tK(0)$. The four-step integration by parts
retains that endpoint term:
\begin{equation}\label{eq:cubic}
 x^4F_{3/2,t}(x)=6tK(0)+\int_0^\infty g^{(4)}(u)\cos(xu)\,du.
\end{equation}
Since $g^{(4)}\in L^1$, the last integral tends to zero by the
Riemann--Lebesgue lemma. This agrees with \eqref{eq:asymptotic}
because $\Gamma(4)=6$ and $\sin(\pi/2)=1$.

One cannot obtain this case by substituting $b=0$ into the preceding
Gamma-cutoff argument: the cutoff $u^{b-1}e^{-u}$ would cease to be
integrable at zero. The separate endpoint calculation is essential.

\section{From the asymptotic to the global zero obstruction}
The coefficient in \eqref{eq:asymptotic} is nonzero: $K(0)>0$,
$\Gamma(2a+1)>0$, and $\sin(\pi(a-1))>0$ for $1<a<2$.
Thus the asymptotic excludes zeros for all sufficiently large positive
real $x$. Evenness excludes them for sufficiently negative $x$.
On the remaining compact real interval, a nonzero entire function
has finitely many zeros. This proves finiteness of the real zero set.

To deduce nonreal infinitude we must also prove that the whole zero
set is infinite. Nonpolynomiality alone is insufficient: exponential
functions are nonpolynomial and zero free. The subquadratic growth
bound supplies the necessary additional rigidity.

\begin{lemma}[Finite-zero rigidity]\label{lem:finitezeros}
Let $f$ be a nonzero entire function satisfying
$|f(z)|\leq C e^{|z|^{3/2}}$ for some $C>0$. If $f$ has finitely
many zeros, then $f(z)=P(z)e^{\alpha z}$ for a nonzero complex
polynomial $P$ and $\alpha\in\C$. Such a function cannot tend to
zero along both real tails.
\end{lemma}
\begin{proof}
Remove every zero with its multiplicity, forming a finite polynomial
$P_0$. Removable singularities make $h=f/P_0$ a zero-free entire
function. A nonzero polynomial has a positive lower bound outside
a sufficiently large disk. On the complementary disk $h$ is bounded
by continuity. Therefore $|h(z)|\leq D e^{|z|^{3/2}}$ for some $D>0$.
This justification for the quotient bound is needed; it does not
follow from a numerator bound without controlling the denominator.

The zero-free entire function $h$ has an entire logarithm $\ell$.
For instance, integrate the entire function $h'/h$ on the simply
connected plane and choose the additive constant at zero. Then
$\operatorname{Re}\ell(z)\leq\log D+|z|^{3/2}$.
The Borel--Carath\'eodory inequality on disks of radii $R$ and $2R$
gives
\[
 \max_{|z|\leq R}|\ell(z)-\ell(0)|=O(1+R^{3/2}).
\]
Cauchy's estimates show for every integer $m\geq2$ that
$|\ell^{(m)}(0)|\leq m!O(1+R^{3/2})/R^m\to0$.
Thus the Taylor series is affine, $\ell(z)=\ell(0)+\alpha z$.
Absorb $e^{\ell(0)}$ into $P_0$ to obtain the claimed factorization.

Choose a sign $s\in\{1,-1\}$ with $s\operatorname{Re}\alpha\geq0$.
Along $z=sx$, $x\to+\infty$, the exponential has modulus at least
one and the nonzero polynomial has a positive lower bound. Hence
$f(sx)$ cannot tend to zero.
\end{proof}

The Borel--Carath\'eodory estimate used here is the disk inequality
\[
 \max_{|z|\leq r}|h(z)-h(0)|
 \leq\frac{2r}{R-r}
 \left(\max_{|z|\leq R}\operatorname{Re}h(z)
       -\operatorname{Re}h(0)\right),\quad 0<r<R,
\]
for holomorphic $h$ near the closed disk. It follows by applying the
maximum principle and Schwarz's lemma to
$(h-h(0))/(2M-(h-h(0)))$, where
$M>\max\operatorname{Re}(h-h(0))$, and then letting $M$ decrease.
This also explains why a bound on the real part suffices for the
Taylor-coefficient argument.

\begin{proof}[Completion of Theorem~\ref{thm:main}]
The density in \eqref{eq:flow} is integrable, so its cosine transform
tends to zero as $x\to\pm\infty$. This form of the Riemann--Lebesgue
lemma follows by approximating the density in $L^1$ by finite step
functions with bounded interval support: their cosine integrals
decay, and the error is bounded by the $L^1$ distance uniformly in $x$.
Lemma~\ref{lem:entire} and Lemma~\ref{lem:finitezeros} now imply
that the total complex zero set is infinite. Removing the finite
real zero set leaves infinitely many nonreal zeros.

If an eventual all-real starting time $T$ existed, choose
$t=\max\{T,1\}>0$. The preceding conclusion at this nonzero time
contradicts the required reality of every zero.
\end{proof}

\begin{corollary}\label{cor:lp}
For $1<a<2$ and $t\neq0$, $F_{a,t}$ is not a locally uniform limit
on $\C$ of real polynomials all of whose roots are real.
\end{corollary}
\begin{proof}
On the upper half-plane such polynomials are zero free, except that
an identically zero polynomial may occur. Local convergence to the
nonzero function $F_{a,t}$ excludes identically zero polynomials
eventually. Hurwitz's theorem then makes the limit zero free there
unless it is identically zero. The same argument applies to the
lower half-plane. Either conclusion contradicts a nonreal zero of
$F_{a,t}$.
\end{proof}

Only the necessary direction of the polynomial-limit criterion is
used. No general converse factorization theorem for the
Laguerre--P\'olya class is needed for this obstruction.

\section{Order convergence on compact sets and on the real axis}
Fix $t$ and $a_0>0$. In a sufficiently small interval around $a_0$,
choose $0<A\leq a\leq B$. For $u\geq1$,
$u^{2a}\leq u^{2B}$, and for $0\leq u\leq1$ it is bounded by one.
For $|z|\leq R$, every fixed spatial derivative of the integrand
is bounded by an integrable multiple of
\[
 (1+u^j)\exp\big(\max(t,0)(1+u^{2B})+Ru\big)K(u).
\]
Continuity of $a\mapsto u^{2a}$ for $u>0$, together with dominated
convergence, therefore gives
\begin{equation}\label{eq:weightedlimit}
 \int_0^\infty u^j e^{Ru}
 |e^{tu^{2a}}-e^{tu^{2a_0}}|K(u)\,du\longrightarrow0.
\end{equation}
At $u=0$ the integrand is defined by continuity; its value on that
single point does not affect the integral.

Taking $j=0$ bounds the supremum of
$|F_{a,t}-F_{a_0,t}|$ on $|z|\leq R$, proving locally uniform
convergence. Taking $j=1$ proves local uniform convergence of the
spatial derivatives. Thus values and derivatives are jointly
continuous in order and spatial argument near any positive order.

For real $x$, $|\cos(xu)|\leq1$, giving the stronger domain bound
\begin{equation}\label{eq:realuniform}
 \sup_{x\in\R}|F_{a,t}(x)-F_{a_0,t}(x)|
 \leq\int_0^\infty
 |e^{tu^{2a}}-e^{tu^{2a_0}}|K(u)\,du\longrightarrow0.
\end{equation}
This is uniform convergence on the entire real axis at fixed time.
It is not uniform convergence on the entire complex plane. The
factor $e^{|\operatorname{Im}z|u}$ would prevent that inference
from \eqref{eq:realuniform}.

These convergence results coexist with Theorem~\ref{thm:main}:
every nearby fractional order can have infinitely many nonreal
zeros, even when a particular limiting transform has none. Neither
small errors on the real axis nor convergence on each fixed disk
controls the complete complex zero set at one common parameter
neighborhood.

\section{Simple real limits force compact escape}
We give the local mechanism explicitly, so the simplicity hypothesis
does not disappear behind a qualitative reference to continuity
of roots.

\begin{lemma}\label{lem:local}
Suppose $a_n\to a_0>0$, $z_n\to w\in\R$,
$F_{a_n,t}(z_n)=0$, and $\operatorname{Im}z_n\neq0$ for every $n$.
Then $F_{a_0,t}(w)=0$ and $F'_{a_0,t}(w)=0$.
\end{lemma}
\begin{proof}
Joint continuity gives the value assertion. Suppose the derivative
at $w$ were $c\neq0$. Joint continuity of the derivative gives a
disk $D$ about $w$ and an order neighborhood such that
\[
 |F'_{a,t}(z)-c|\leq|c|/2\qquad(z\in D)
\]
for all those orders. Since $D$ is convex, integration along the
segment between $z,v\in D$ shows
\[
 |F_{a,t}(z)-F_{a,t}(v)-c(z-v)|
 \leq\frac{|c|}{2}|z-v|.
\]
Hence $|F_{a,t}(z)-F_{a,t}(v)|\geq |c||z-v|/2$.
For all large $n$, both $z_n$ and $\bar z_n$ lie in $D$.
Conjugation symmetry makes both zeros of $F_{a_n,t}$.
The lower bound forces $z_n=\bar z_n$, a contradiction.
\end{proof}

\begin{proposition}\label{prop:compact}
If every zero of $F_{a_0,t}$ is real and simple, then for each
compact set $Q\subset\C$ there is an order neighborhood of $a_0$
in which every zero belonging to $Q$ is real.
\end{proposition}
\begin{proof}
Otherwise select $a_n\to a_0$ and nonreal zeros $z_n\in Q$.
Compactness supplies a subsequence converging to $w\in Q$.
Joint continuity makes $w$ a limiting zero; the hypothesis makes
it real. Lemma~\ref{lem:local} makes it multiple, contrary to simplicity.
\end{proof}

\begin{proof}[Proof of the escape assertions in Theorem~\ref{thm:limit}]
Apply Proposition~\ref{prop:compact} to each closed disk
$\{|z|\leq R\}$. A sequence of nonreal zeros with $a_n\to a_0$
eventually avoids that disk. This is exactly $|z_n|\to\infty$.

For the unconditional selection near one, put
$a_n=1+1/(n+3)\in(1,3/2)$. Theorem~\ref{thm:main} gives an infinite
nonreal zero set for each $a_n$. The nonzero entire function has
only finitely many zeros in any compact disk, so choose one with
$|z_n|>n$. This chosen sequence escapes and needs no hypothesis
about $F_{1,t}$.
\end{proof}

The simplicity assumption cannot be dropped in the general compact
stability principle. For instance, $z^2+\varepsilon$ converges
locally uniformly to $z^2$ as $\varepsilon\downarrow0$, while its
nonreal zeros converge to the double real zero at zero. This
elementary comparison is explanatory; it is not a claim that this
polynomial family is an arithmetic theta flow.

\section{Endpoints, density analyticity, and parameter limits}
At $t=0$, every order in \eqref{eq:flow} gives the same transform
$\int K(u)\cos(zu)\,du$. The coefficient $L(a,0)$ vanishes,
so the preceding nonreal-zero obstruction makes no assertion at
that time. In particular no conclusion about the Riemann hypothesis
is extracted from a theorem whose hypothesis is $t\neq0$.

Likewise $L(1,t)=L(2,t)=0$. Direct differentiation of the explicit
coefficient gives
\begin{equation}\label{eq:coefficientlimit}
 \frac{L(a,t)}{a-1}\longrightarrow2\pi tK(0)
 \qquad(a\downarrow1).
\end{equation}
Indeed, the sine factor has derivative $\pi$ at one and
$\Gamma(2\cdot1+1)=2$; derivatives of the other factors multiply
the vanishing sine. This is a coefficient limit at fixed time,
not a uniform relative asymptotic on the rectangle
$a\downarrow1$, $x\to\infty$.

The density distinction is also exact. At $a=3/2$ its even full-line
form is $e^{t|u|^3}K(u)$. For $t\neq0$ the right and left third
derivatives at zero differ by $12tK(0)$, so this full-line density
cannot be the restriction of a holomorphic germ through zero.
Its positive-axis form $e^{tu^3}K(u)$ does have a holomorphic
extension there. These are different functions on the negative
axis. Nevertheless its Fourier transform is entire by
Lemma~\ref{lem:entire}. Entire dependence on frequency follows
from decay and does not require analyticity of the full-line density
at its endpoint.

The same distinction prevents an unjustified complex branch choice
for a nonintegral real exponent. Definition \eqref{eq:flow} uses
the real power of a nonnegative real variable. The proof of its
entireness differentiates in $z$, not in a fictitious complex
continuation of $u^{2a}$ across zero.

The library separately proves eventual real and simple zeros at
quartic order. Combined with the compact escape principle, that
result gives an application with $a_n=2-1/(n+3)$. Its global
quartic saddle analysis is a separate theorem, not reconstructed
in this paper and not an assumption in Theorem~\ref{thm:main}.
No effective upper time, nonreal-zero escape rate, or classification
for real orders above two follows from the present argument.



\section{Uniform absolute errors and a real-frequency transition}
The fixed-order cancellation estimates of Section 3 do not alone give
uniformity as $a\downarrow1$. The coefficient of their singular
derivative must be retained. This leads to a quantitative real-frequency
statement, while still excluding relative uniformity on a full rectangle.

\begin{theorem}\label{thm:transition}
For each fixed real $t$, there is $A_t\geq0$ such that
\begin{equation}\label{eq:uniformabsolute}
 \bigl|x^{2a+1}F_{a,t}(x)-L(a,t)\bigr|\leq A_t x^{-1/2}
 \quad\left(1<a\leq\frac54,\ x\geq1\right).
\end{equation}
If $t\ne0$, there are $0<\delta_t\leq1/4$, $C_t\geq0$ and
$K_t>0$ such that, for $1<a<1+\delta_t$,
\begin{align}
 \left|\frac{x^{2a+1}F_{a,t}(x)}{L(a,t)}-1\right|
 &\leq\frac{C_t}{(a-1)\sqrt{x}}\quad(x\geq1),
       \label{eq:uniformrelative}\\
 F_{a,t}(x)=0,\ x\in\R\quad&\Longrightarrow\quad
 |x|\leq\max\left\{1,\left(\frac{K_t}{a-1}\right)^2\right\}.
       \label{eq:realzeroradius}
\end{align}
In particular the relative asymptotic holds along every sequence
$a_n>1$, $a_n\to1$, $x_n\to\infty$ with
$(a_n-1)\sqrt{x_n}\to\infty$.
\end{theorem}


\subsection*{A surviving boundary mass behind pointwise vanishing}
There is an exact measure-level reason to keep the small coefficient
in the cutoff. For $a>1$, put $b=2a-2>0$ and define the real
signed density on $u>0$
\[
 c_{a,t}(u)=tK(0)(2a)(2a-1)(2a-2)u^{2a-3}e^{-u},
 \qquad m(a,t)=tK(0)\Gamma(2a+1).
\]
This is precisely the real value of the cutoff subtracted from the
third derivative, not the full deformed arithmetic density.
The Gamma integral and its recurrence give
\begin{align}\label{eq:cutoffmassmoments}
 \int_0^\infty c_{a,t}(u)\,du&=m(a,t),&
 \int_0^\infty|c_{a,t}(u)|\,du&=|m(a,t)|,\\
 \int_0^\infty u|c_{a,t}(u)|\,du&=b|m(a,t)|.\notag
\end{align}
Indeed the integrals of $u^{b-1}e^{-u}$ and $u^be^{-u}$
are $\Gamma(b)$ and $\Gamma(b+1)=b\Gamma(b)$,
respectively; three recurrences yield
$b(b+1)(b+2)\Gamma(b)=\Gamma(b+3)=\Gamma(2a+1)$.
All coefficients other than $t$ are positive. This proves both
the signed and absolute formulas, as well as integrability and
the first absolute moment.
\par\noindent\code{RiemannRealOrderCutoffBoundary.cutoff_integral}.
\par\noindent\code{RiemannRealOrderCutoffBoundary.cutoff_norm_integral}.
\par\noindent\code{RiemannRealOrderCutoffBoundary.cutoff_norm_first_moment}.

If $f:\R\to\C$ is globally Lipschitz with constant $L$, then
$|f(u)|\leq|f(0)|+Lu$ on the positive half-line, so the
preceding moments justify every integral in
\begin{equation}\label{eq:cutofftestbound}
 \left|\int_0^\infty c_{a,t}(u)f(u)\,du-m(a,t)f(0)\right|
 \leq Lb|m(a,t)|.
\end{equation}
To prove it, subtract $f(0)$ under the integral, use
$|f(u)-f(0)|\leq Lu$, and apply the last identity in
\eqref{eq:cutoffmassmoments}. Since Gamma is continuous at $3$,
$m(a,t)\to2tK(0)$ as $a\downarrow1$; hence
\begin{equation}\label{eq:cutoffweak}
 \int_0^\infty c_{a,t}(u)f(u)\,du\longrightarrow2tK(0)f(0).
\end{equation}
This holds also for unbounded Lipschitz tests, and for either sign
of $t$. At $t=0$ the density and all the displayed quantities vanish.
\par\noindent\code{RiemannRealOrderCutoffConcentration.cutoff_lipschitz_error_bound}.
\par\noindent\code{RiemannRealOrderCutoffConcentration.cutoff_lipschitz_tendsto}.

For $t\ne0$ the normalized density is exactly the positive
Gamma probability density
\[
 \frac{c_{a,t}(u)}{m(a,t)}
       =\frac{u^{b-1}e^{-u}}{\Gamma(b)}.
\]
Let $\nu_a$ denote this probability on $[0,\infty)$, with no
atom at zero. Its first moment is $b$. The preceding calculation
therefore gives the exact human consequence
\begin{equation}\label{eq:cutofflipmetric}
 \sup_{\operatorname{Lip}(f)\leq1}
 \left|\int f\,d\nu_a-f(0)\right|=b,
\end{equation}
where tests are real functions on $[0,\infty)$ with no boundedness
restriction. The upper bound is the first-moment inequality; the
test $f(u)=u$ attains it. Integrals exist because such tests grow
at most linearly. This is an explicit Gamma approximate identity,
not a new probability law associated with the entire transform.

It also converges against every bounded continuous test, a human
consequence beyond the named globally Lipschitz formulation.
For $\eta>0$, the first moment implies
$\nu_a([\eta,\infty))\leq b/\eta$. Splitting the test error
at $\eta$ gives
\[
 \left|\int f\,d\nu_a-f(0)\right|
 \leq\sup_{0\leq u\leq\eta}|f(u)-f(0)|
           +2\|f\|_\infty b/\eta.
\]
First let $a\downarrow1$ at fixed $\eta$, then let
$\eta\downarrow0$. Thus the signed cutoff measures converge
weakly on the \emph{closed} half-line to $2tK(0)\delta_0$.
The endpoint is included in this statement; the original density
integrals themselves remain over $u>0$.

Nevertheless $c_{a,t}(u)\to0$ for every fixed $u>0$: the factor
$2a-2$ tends to zero and the other factors have finite limits there.
For $t\ne0$, \eqref{eq:cutoffmassmoments} instead gives
$\int|c_{a,t}|\to|2tK(0)|>0$, because $K(0)>0$.
There is no common integrable dominator on any order neighborhood
$1<a<1+\delta$. Otherwise take the sequence
$a_n=1+\delta/(n+2)$; the countable intersection of its
almost-everywhere domination events and dominated convergence
would force its integrals to zero, contradicting the mass limit.
This is precisely why pointwise disappearance cannot remove this
cutoff from the parameter-limit analysis. It does not contradict
the already proved uniform convergence of the original transforms,
nor does it alter any zero classification.
\par\noindent\code{RiemannRealOrderCutoffConcentration.cutoff_pointwise_tendsto_zero}.
\par\noindent\code{RiemannRealOrderCutoffConcentration.no_common_integrable_majorant}.

\subsection*{A uniform integrable derivative after cancellation}
Fix $t$, put $\beta=2a$, $b=\beta-2\in(0,1/2]$ and
$E(u)=e^{tu^\beta}$. Keep $R_3$ and $C_3$ from
\eqref{eq:remainder}. Write
\[
 p_j(u)=t(\beta)_j u^{\beta-j},\qquad
 D(u)=E(u)K(u)-K(0)e^{-u}.
\]
All constants in this subsection may depend on $t,K$, but are
chosen before $a$. On $0<u\leq1$, $E,p_1,p_2$ are uniformly
bounded. The derivative bound on $K$ and the mean-value estimates
$|e^v-1|\leq |v|e^{|v|}$ and $1-e^{-u}\leq u$ give
$|D(u)|\leq C u$: use
\[
 D=E(K-K(0))+(E-1)K(0)+K(0)(1-e^{-u}),\quad
 u^\beta\leq u.
\]
Since $2<\beta\leq5/2$,
\begin{equation}\label{eq:retainedb}
 |p_3(u)|\leq Cb u^{b-1},\qquad
 |p_4(u)|\leq Cb u^{b-2},\qquad
 \int_0^1 b u^{b-1}\,du=1.
\end{equation}
The last equality follows by evaluating $u^b$ at the endpoints;
the small factor $b$ must not be discarded.

For clarity the entire fourth-derivative polynomial is
\begin{align*}
 g^{(4)}=E\bigl[&K^{(4)}+4p_1K^{(3)}+6(p_2+p_1^2)K''
       +4(p_3+3p_1p_2+p_1^3)K'\\
       &+(p_4+4p_1p_3+3p_2^2+6p_1^2p_2+p_1^4)K\bigr].
\end{align*}
It follows by repeatedly differentiating $g=EK$, using $E'=p_1E$
and $p_j'=p_{j+1}$. Remove just the term $Ep_4K$ and call
the remaining expression $H(u)$. Every remaining occurrence of
$p_3$ is linear with uniformly bounded coefficient. All other
factors are uniformly bounded on $(0,1]$, so
$|H(u)|\leq C(1+|p_3(u)|)$. Differentiating the actual cutoff
$C_3u^{\beta-3}e^{-u}=p_3K(0)e^{-u}$ gives the exact cancellation
\[
 R_3'(u)=p_4(u)D(u)+H(u)+p_3(u)K(0)e^{-u}.
\]
Thus \eqref{eq:retainedb} bounds its norm by
$C(1+b u^{b-1})$ on $(0,1]$, whose integral is bounded uniformly
in $b$. No common pointwise integrable majorant for
$b u^{b-1}$ is asserted or needed.

On $u\geq1$, the finitely many phase derivatives have a common
polynomial bound in $u$, and $E(u)\leq e^{|t|u^{5/2}}$.
Lemma~\ref{lem:kernel} bounds every $K^{(j)}$, $j\leq4$, by
$C_j\exp(C_ju-c e^{4u})$ for $c>0$. Multiplication by the
polynomial phase factors therefore leaves a common integrable
bound, since $e^{4u}$ dominates $u^{5/2}$ and every linear or
logarithmic term. The cutoff derivative is bounded by $Ce^{-u}$:
its powers are $u^{b-1},u^{b-2}\leq1$ and its coefficients are
uniformly bounded. Combining the two regions proves
\begin{equation}\label{eq:uniformRprime}
 \sup_{1<a\leq5/4}\int_0^\infty|R_3'(u)|\,du<\infty.
\end{equation}
For each order the endpoint and infinity values of $R_3$ vanish,
as proved in Lemma~\ref{lem:cancel}. One integration by parts
therefore yields
\[
 \left|\int_0^\infty R_3(u)\sin(xu)\,du\right|\leq\frac Mx.
\]
Multiplying by $x^b$ gives at most
$Mx^{b-1}\leq Mx^{-1/2}$ for $x\geq1$.
This is the classical derivative-variation error mechanism, also
recorded in DLMF \cite[\S2.3(i)]{DLMF}; the coefficient-aware
endpoint estimate above is what makes its constant uniform here.

\subsection*{The cutoff transform has a uniform reciprocal-frequency rate}
Let $m_t(\beta)=tK(0)\Gamma(\beta+1)$, continuous and bounded
for $2\leq\beta\leq5/2$. Only this coefficient and the power profile are extended to $b=0$ for bounding purposes; the actual cutoff integral below still has $b>0$. Equation~\eqref{eq:gamma} and the
Gamma recurrence give, for $x>0$,
\[
 x^b C_3\int_0^\infty u^{b-1}e^{-u}e^{ixu}\,du
      =m_t(\beta)B(1/x)^b,\qquad B(y)=(y-i)^{-1}.
\]
The power is the principal one. On $0\leq y\leq1$, $B(y)$
stays in the open upper half-plane, $1/\sqrt2\leq|B(y)|\leq1$
and $|B'(y)|\leq1$. For $0\leq b\leq1/2$, differentiation
gives
\[
 \left|\frac{d}{dy}B(y)^b\right|
      =|b B(y)^{b-1}B'(y)|\leq\frac{\sqrt2}{2}.
\]
The branch is valid throughout the segment, so integration of this
derivative gives $|B(1/x)^b-i^b|\leq C/x$ uniformly in $b$
for $x\geq1$. Taking imaginary parts, the limit value is
\[
 m_t(\beta)\sin(\pi b/2)=L(a,t).
\]
The cutoff error is therefore at most $Q_t/x$ uniformly in order.
The identity from Section 3 is
\[
 x^3F_{a,t}(x)=\int_0^\infty g^{(3)}(u)\sin(xu)\,du.
\]
Together with \eqref{eq:uniformRprime}, it proves
\eqref{eq:uniformabsolute}, since $x^{-1}\leq x^{-1/2}$.

\subsection*{Division by the vanishing coefficient and its limitation}
For $t\ne0$, \eqref{eq:coefficientlimit} gives constants
$c_t>0$ and $0<\delta_t\leq1/4$ with
$|L(a,t)|\geq c_t(a-1)$ for $1<a<1+\delta_t$.
Divide \eqref{eq:uniformabsolute} by this positive lower bound to
obtain \eqref{eq:uniformrelative} with $C_t=A_t/c_t$.
If $(a-1)\sqrt{x}>K_t:=C_t+1$, the right side is less than one,
whereas a zero of the flow would make the left side equal one.
Evenness supplies the same conclusion for negative real $x$.
Squaring the resulting bound on $\sqrt{|x|}$, and treating
$|x|<1$ separately, proves \eqref{eq:realzeroradius}.
The diagonal relative limit follows directly from the bound.
The constants exist for each fixed $t$; no numerical values,
optimal power, complex-zero radius or uniformity in $t$ is claimed.

The full rectangular relative assertion is actually false. For any
$\delta>0$, real $R$ and finite $M$, there are $1<a<1+\delta$
with $a<3/2$ and $x\geq\max(R,1)$ such that
\begin{equation}\label{eq:rectangularfailure}
 \left|\frac{x^{2a+1}F_{a,t}(x)}{L(a,t)}-1\right|>M
                       \qquad(t\ne0).
\end{equation}
Indeed $F_{1,t}$ is a nonzero entire function: its value at zero
is the positive integral of $e^{tu^2}K(u)$. Its real zeros are
isolated, so choose such an $x$ with $F_{1,t}(x)\ne0$.
At that fixed $x$, continuity in order gives
$x^{2a+1}F_{a,t}(x)\to x^3F_{1,t}(x)\ne0$, while
$L(a,t)\to0$ and is nonzero for $1<a<3/2$. The absolute
relative error therefore tends to infinity as $a\downarrow1$,
proving \eqref{eq:rectangularfailure}. This explains why a
parameter-dependent region is needed; it does not assert that the
particular transition scale in \eqref{eq:uniformrelative} is sharp.

\section{Gaussian translation needs control beyond compact sets}
The compact convergence proved above has a precise use in local zero
arguments. It does not, by itself, justify passage through a Gaussian
average along an unbounded imaginary line. We record a positive identity
for the actual quadratic flow and a separate polynomial counterexample
to that proposed continuity inference.

For $s>0$ let $\nu_s=N(0,2s)$, with density
$(4\pi s)^{-1/2}e^{-y^2/(4s)}$; put $\nu_0=\delta_0$.
Completing the Gaussian square gives
\[
 \int_{\R} e^{by}\,d\nu_s(y)=e^{sb^2}\quad(b\in\R).
\]
The complex version is $\int e^{qy}\,d\nu_s=e^{sq^2}$.
For completeness, this follows for real $q$ by the same square
completion; both sides are entire in $q$. Differentiation under the
integral is justified on $|q|\leq R$ by the integrable bounds
$|y|^j e^{R|y|}$ against the Gaussian. The identity theorem supplies
the complex equality. At $s=0$ it is immediate.

\begin{proposition}\label{prop:gaussiantranslate}
For every real $t$, $s\geq0$ and $z\in\C$, the following are
genuine absolutely integrable Gaussian averages:
\begin{equation}\label{eq:actualgaussian}
 \int F_{1,t}(z+iy)\,d\nu_s(y)=F_{1,t+s}(z),\qquad
 \int F_{1,t}(z+y)\,d\nu_s(y)=F_{1,t-s}(z).
\end{equation}
\end{proposition}
\begin{proof}
For the imaginary shift, the absolute integrand in the double
cosine integral is bounded by
\[
 e^{tu^2}K(u)e^{|\operatorname{Im}z|u}e^{u|y|}.
\]
Since $e^{u|y|}\leq e^{uy}+e^{-uy}$, its integral in $y$ is at
most $2e^{(t+s)u^2+|\operatorname{Im}z|u}K(u)$. This is integrable
in $u\geq0$ by Lemma~\ref{lem:kernel}; in particular product
integrability is proved before Fubini is used. Expressing the cosine
as the mean of two complex exponentials and applying the Gaussian
formula multiplies it by $e^{su^2}$, giving the first identity.
For the real shift, the absolute bound is already
$e^{tu^2+|\operatorname{Im}z|u}K(u)$, independent of $y$.
The Gaussian characteristic formula multiplies the cosine by
$e^{-su^2}$, giving the second identity. At $s=0$ both reduce
to evaluation.
\end{proof}
These are classical Gaussian translation identities with the actual
arithmetic integrability checked. Their time signs differ because
the displacement is imaginary or real. They give no conclusion
about quadratic zero reality without a further real-zero theorem.

Define the time-one imaginary Gaussian functional on integrable
entire functions by
\[
 \mathcal H f=\int_{\R}f(iy)\,d\nu_1(y).
\]
For $m\geq1$ put
\begin{equation}\label{eq:heatpolynomials}
 q_m(z)=(-1)^m\frac{m!}{(2m)!}z^{2m},\qquad p_m(z)=1+q_m(z).
\end{equation}
These are even polynomials with real coefficients and $p_m(0)=1$.
Since $\binom{2m}{m}\geq1$, one has $(m!)^2\leq(2m)!$.
Thus, on $|z|\leq R$,
\[
 |q_m(z)|\leq\frac{R^{2m}}{m!}\longrightarrow0.
\]
The convergence follows, for example, because these are terms of
the convergent exponential series $\sum_mR^{2m}/m!$.
Consequently $p_m\to1$ uniformly on every compact subset of $\C$.
The same estimate also gives the common, but overly large, bound
$|p_m(z)|\leq1+e^{|z|^2}$.

Gaussian integration by parts with density
$e^{-y^2/4}/\sqrt{4\pi}$ yields
\[
 \int y^{2m}\,d\nu_1(y)=2(2m-1)\int y^{2m-2}\,d\nu_1(y)
                         =\frac{(2m)!}{m!}.
\]
The boundary terms vanish by polynomial Gaussian decay; iteration
starts with mass one. On the imaginary axis the two signs cancel:
\[
 q_m(iy)=\frac{m!}{(2m)!}y^{2m}\geq0.
\]
Every polynomial is absolutely integrable there, but nevertheless
\begin{equation}\label{eq:heatdiscontinuity}
 \mathcal Hq_m=1,\qquad \mathcal Hp_m=2,\qquad
 \mathcal H1=1,\qquad \int |p_m(iy)-1|\,d\nu_1(y)=1.
\end{equation}
Hence $\mathcal H$ is not sequentially continuous for compact-open
convergence, even on normalized even real polynomials. At the
polynomial level the same value is obtained by the finite differential
expression
\[
 (e^{-\partial_z^2}p_m)(0)
 :=\sum_{j=0}^{m}\frac{(-1)^j}{j!}p_m^{(2j)}(0)=2:
\]
only $j=0,m$ contribute. This equality is a finite algebraic
calculation, not a definition of that exponential on arbitrary entire
functions.

There is no common $g\in L^1(\nu_1)$ dominating $|p_m(iy)|$
almost everywhere for every $m$. Otherwise countability permits a
common exceptional-null-set removal and dominated convergence would
force the integrals to tend to one, contrary to
\eqref{eq:heatdiscontinuity}. A further elementary human consequence
locates the missing mass: for every fixed $R$, compact convergence
gives $\int_{|y|\leq R}q_m(iy)\,d\nu_1\to0$, whereas the total
integral is one. Its complementary integral therefore tends to one.
Pointwise convergence has lost uniform integrability at infinity.

For comparison, a sufficient condition is a common bound
$|f_m(iy)|\leq C e^{\tau y^2}$ with $\tau<1/4$, together with
compact convergence of $f_m$ to $f$. The same bound holds for the
limit, and its product with the Gaussian density is a constant
multiple of $e^{-(1/4-\tau)y^2}$, which is integrable. Dominated
convergence then does give $\mathcal H f_m\to\mathcal H f$.
This is a labelled classical sufficient condition, not a necessary
growth classification or an additional named formal theorem.
The polynomials \eqref{eq:heatpolynomials} have not been asserted
real-rooted. They therefore do not refute any theorem with an
additional split-polynomial or growth hypothesis. The example
isolates the missing topology assumption; it changes none of the
actual order-limit or zero conclusions proved in this paper.

\section{Formal correspondence and editorial status}
All principal results are supported by the frozen commit
\begin{center}\small\ttfamily
bab688bb6140cd6316c80680095a914ea0d346db.
\end{center}
The namespace prefix used below is
\code{OperatorDerivativeRewrite.Foundation}.
Exact source links are in \code{theorem-map.json}, together with
line numbers and SHA256 hashes at this commit.

\begin{longtable}{@{}p{0.25\textwidth}p{0.69\textwidth}@{}}
\textit{Paper content}&\textit{Lean namespace and declaration}\\\hline
Kernel, actual integral&\code{RiemannQuarticKernel.kernel};
\code{RiemannRealOrderFlow.flow_eq_cosine_integral}\\[3pt]
Lemma~\ref{lem:entire}&\code{RiemannRealOrderFlow.flow_differentiable},
\code{actual_flow_zero_pos};
\code{RiemannRealOrderComplexZeros.actual_subquadratic_growth}\\[3pt]
First fractional range&\code{RiemannRealOrderFractionalCoefficient.actual_flow_scaled_limit}\\[3pt]
Second fractional range&\code{RiemannRealOrderFourthAsymptotic.actual_flow_scaled_limit}\\[3pt]
Cubic endpoint&\code{RiemannRealOrderCubicEndpoint.actual_flow_scaled_limit}\\[3pt]
Lemma~\ref{lem:finitezeros}&\code{SubquadraticFiniteZeroFactorization.finite_zero_exponential_polynomial},
\code{exponential_polynomial_not_real_decay}\\[3pt]
Theorem~\ref{thm:main}&\code{RiemannRealOrderQuarticLimitZeros.between_one_two_real_zeros_finite},
\code{between_one_two_nonreal_zeros_infinite},
\code{between_one_two_no_eventual_real_zeros}\\[3pt]
Corollary~\ref{cor:lp}&\code{RiemannRealOrderQuarticLimitZeros.between_one_two_no_split_polynomial_limit}\\[3pt]
Order convergence&\code{RiemannRealOrderLimit.flow_order_tendsto_locally_uniformly};
\code{RiemannRealOrderRealUniform.flow_order_tendsto_uniformly_real}\\[3pt]
Lemma~\ref{lem:local}, Proposition~\ref{prop:compact}&\code{RiemannRealOrderZeroLimits.derivative_zero_of_nonreal_order_limit},
\code{eventually_real_zeros_on_compact}\\[3pt]
Theorem~\ref{thm:limit}&\code{RiemannRealOrderZeroLimits.nonreal_zero_norm_tendsto_atTop};
\code{RiemannRealOrderFractionalZeros.exists_approaching_escaping_nonreal_zeros}\\[3pt]
\eqref{eq:coefficientlimit}&\code{RiemannRealOrderFractionalCoefficient.leading_div_order_gap_limit}\\
\end{longtable}

Section 8 is mapped to the uniform endpoint, regular-part,
tail and cutoff bounds, the actual transition-rate theorems and
the rectangular relative obstruction. The explicit Bell-polynomial
and reciprocal-profile estimates are human reconstructions of that
verified mechanism; their constants are not claimed to be optimized.
The new uniform Fourier and zero-radius conclusions rest on the full
analytic proofs, not a numerical frequency scan.

Section 9 uses the actual Gaussian averages in \code{RiemannQuadraticGaussian} and the polynomial, integrability and nonconvergence statements in \code{GaussianHeatTopology}. Their exact declarations are in the theorem map. The compact-tail mass and sufficient Gaussian-growth criterion are fully proved classical human deductions. Thirty-six exact finite checks verify moment, sign and normalization transcriptions; they do not establish the analytic limits.

The elementary pointwise estimates in the kernel and canceled-remainder
proofs are human reconstructions of the verified mechanism. They
are not claimed as additional named Lean declarations with the same
big-$O$ notation. The actual derivatives, integrability, endpoint
limits, integration-by-parts identities and final convergence
statements have their own exact declarations. In particular, the
fourth-order cancellation uses a separately verified actual fifth
derivative; a formal exponential jet defined only through order four
is not differentiated as if it were that fifth derivative.
The density obstruction is separately recorded as
\code{RiemannRealOrderBranchObstruction.three_halves_actual_density_obstruction}.

The baseline audit records a complete successful root build and
54,943 declarations audited with only \code{propext},
\code{Classical.choice}, and \code{Quot.sound}. The editorial pass
does not rerun a redundant full Lean build or change mathematical
sources. Hash-linked provenance, compilation logs, and visual-review
receipts accompany this draft.

P\'olya's original 1927 discussion~\cite[Chapter I, no. 6,
assertion II, pp. 14--15]{polyaJensen} already proves the smooth
odd-endpoint obstruction: infinitely many zeros in a fixed horizontal
strip force every odd endpoint derivative to vanish. For
$e^{tu^3}K(u)$ the third derivative is $6tK(0)\neq0$.
Thus finiteness of its real zeros is a classical consequence; the
coefficient in~\eqref{eq:cubic} is made explicit here. The nonintegral
densities do not satisfy that statement's smoothness premise.

The contribution is the precise arithmetic statement and its
readable proof, with exact endpoint and limit safeguards. General
endpoint Fourier asymptotics, finite-zero rigidity, Hurwitz's theorem,
and the role of simple zeros are established analysis. Priority for
the arithmetic specialization remains unresolved. A publication
decision should compare it directly with P\'olya-type non-even-power
transforms and subsequent real-zero-preserving transform criteria.

\clearpage
\begin{thebibliography}{9}
\bibitem{polya} G. P\'olya,
On the zeros of an integral function represented by Fourier's integral,
\emph{Messenger of Mathematics} \textbf{52} (1923), 185--188.
Primary paper identified in the preliminary literature search;
detailed historical theorem comparison remains pending.
\bibitem{polyaJensen} G. P\'olya,
\emph{\"Uber die algebraisch-funktionentheoretischen Untersuchungen
von J. L. W. V. Jensen},
Det Kgl. Danske Videnskabernes Selskab,
Mathematisk-fysiske Meddelelser \textbf{7}, no. 17 (1927), 3--33.
\href{https://gymarkiv.sdu.dk/MFM/kdvs/mfm\%201-9/mfm-7-17.pdf}{Primary text}.
\bibitem{cardon} D. A. Cardon,
Fourier transforms having only real zeros,
\emph{Proc. Amer. Math. Soc.} \textbf{133} (2005), 1349--1356
(published electronically October 18, 2004).
\href{https://doi.org/10.1090/S0002-9939-04-07677-4}{doi:10.1090/S0002-9939-04-07677-4}.
\bibitem{DLMF} NIST Digital Library of Mathematical Functions,
\emph{Integrals of a real variable}, \S2.3(i)--(ii).
\href{https://dlmf.nist.gov/2.3}{Official reference}.
Derivative-variation and endpoint-method context only.
\end{thebibliography}
\end{document}
